For routed universe $U(q)$, JHQ's refinement multiplier sets \cite{jhq}
\begin{equation}
 c(\alpha)=\max\{k,\lceil\alpha k\rceil\},\qquad
 C_\alpha(q)=\operatorname{Top}^{\rm primary}_{\min\{c(\alpha),|U(q)|\}}U(q).
 \label{eq:6-budget}
\end{equation}
$C_\alpha$ contains valid survivors; underfull output slots receive sentinels. Analysis requires at least $k$ valid candidates. $\alpha$ changes retention within $U$, while probes change $U$ itself.

Too much refinement wastes residual reads; too little can lose neighbours. The grid spans 20--2,000 survivors at $k=10$ without retraining. Smaller budgets do not release stored residuals or necessarily shrink previously grown workspace.

\begin{figure}[tbp]
\centering
% Logical work boundary, not a measured latency or pruning diagram.
\begin{tikzpicture}[
 font=\footnotesize\linespread{1}\selectfont,
 box/.style={draw=ThesisNavy!65,rounded corners=2pt,align=center,
 text width=2.15cm,minimum height=1.5cm,inner sep=3pt},
 arr/.style={-{Latex[length=2mm]},draw=ThesisNavy!80,thick}]
\node[box] (route) at (0,0) {Probe IVF lists\\Fixed $n_{\rm probe}$};
\node[box] (primary) at (3,0) {Score $U(q)$\\Primary codes};
\node[box,fill=ThesisNavy!7] (select) at (6,0)
 {Primary top-$c$\\Keep $C_\alpha(q)$};
\node[box,fill=ThesisNavy!7] (refine) at (9,0)
 {Refine $c(\alpha)$ slots\\Residual rows};
\node[box,fill=ThesisNavy!7] (final) at (12,0)
 {Final top-$k$\\Return $k$ IDs};
\draw[arr] (route) -- (primary);
\draw[arr] (primary) -- (select);
\draw[arr] (select) -- (refine);
\draw[arr] (refine) -- (final);
\node[align=center,text width=5.1cm] at (1.5,1.4)
 {Fixed routed universe and number of primary scores};
\node[align=center,text width=7.6cm] at (9,1.4)
 {Budget-dependent selection, residual scoring and final selection};
\draw[densely dashed,ThesisNavy!45] (4.5,-0.9) -- (4.5,1.1);
\node[align=center,text width=13.8cm] at (6,-1.45)
 {$c(\alpha)=\max\{k,\lceil\alpha k\rceil\}$ slots;
 valid survivors are capped by $|U(q)|$.\\
 The stored model and residual codes are reused across budgets.};
\end{tikzpicture}

\caption[Work affected by the refinement budget]{The refinement budget changes work after primary-score
production. With a fixed query, index and probe count, it leaves the
routed candidate universe unchanged. The highlighted stages depend on
survivor capacity; this is a logical work decomposition, not a measured
latency breakdown. The scan kernel also performs budget-dependent
admission and compaction. Padding applies when fewer than $c$ candidates exist.}
\label{fig:6-budget-scope}
\end{figure}

For fixed batch and routing, the cost split is
\begin{equation}
 T(\alpha)=T_{\rm base}+T_{\rm select}(c(\alpha))+
 T_{\rm refine}(c(\alpha))+T_{\rm final}(c(\alpha)).
 \label{eq:6-cost-model}
\end{equation}
$T_{\rm base}$ covers transfers, rotation, routing and primary scores; survivor-dependent work follows \cref{fig:6-budget-scope}. Capacity/resource thresholds prevent a monotone-latency claim. The rule minimises accepted budget, not noisy QPS (\cref{sec:7-budget-trace}).


Each calibration trial still runs rotation, probing and primary scoring
before refinement. Reusing the index avoids rebuilding or retraining;
it does not cache every intermediate query result across budget trials.
The cost of calibration must therefore count full searches, not merely
the difference between two residual-scoring workloads.

The policy selects one budget for an index, routing configuration, $k$
and query workload. It uses neither a trained per-query predictor nor
an online choice for each arriving query. Scan-group scheduling remains
the separate mechanism of \cref{sec:5-5}; its execution path must still
be considered when interpreting budget calibration.

\paragraph{Position among adaptive ANN methods.}
$\alpha$ serves the candidate-expansion role of Faiss's \texttt{k\_factor} and cuVS refinement \cite{faiss_refine,cuvs_refinement}. Alternatives use GT tuning \cite{faiss_autotune}, learned per-query stopping \cite{li_termination} or predicted recall \cite{darth}. Here unlabelled agreement chooses one reusable workload budget. Scores still use reconstructed vectors, unlike the cuVS source-vector refinement example.